\scanpage{5}
\(\overset{\hspace{0.55em}\unsure{?|}}{\text{\uline{NB}}}\)\quad Si $\sigma$ est un $\underset{\hspace{-1.3em}\unsure{\cdot|}}{\text{automorphisme}}$ du corps des complexes, et si $f$ est ``spécial'', alors $f^\sigma$ est \unsure{égalmt|également} spécial -- donc $\operatorname{Aut}(\mathbb C)$ opere sur l'ens des $f$ spèciaux. Mais il opère par l'intermédiaire de $\operatorname{Aut}(\unsure{\overline{\mathrm Q}|\overline{\mathbb Q}})$, car si $f=P/Q$ spécial, les $\underset{\hspace{1.5em}\unsure{-|=|}}{\text{coeff.}}$ de $P,Q$ sont des \uline{nbs algébriques}. Il revient au $=$ de dire que les zéros de $P,Q$ sont des nbs algébriques, i.e. le diviseur de $f$ est

% A pale hook and dark dot above NB are retained as an unsure question-like shape.
% A tiny detached dot below automorphisme is retained; its function is uncertain.
% Faint detached strokes below the right end of coeff. are preserved above;
% the mark's origin and intention remain unresolved (full scan x409--436, y887--902).
% The following drawing preserves the relative placement of four detached
% lower mark-groups. The three uncertain groups have unresolved intention;
% no mathematical relation between them is supplied.
\figure{\begin{tikzpicture}[x=0.011cm,y=0.011cm]
  \node[inner sep=0pt] at (489,161) {$\unsure{\cdot\,\cdot|-}$};
  \node[inner sep=0pt] at (0,0) {$\unsure{-|=}$};
  \node[inner sep=0pt] at (61,5) {$\unsure{\frown|-|}$};
  \node[inner sep=0pt] at (545,0) {$-$};
\end{tikzpicture}}
% archive sheet number: 4
